% !TEX root = ../monograph.tex
\chapter[Reflection spaces and Ward identities]{Reflection spaces, Ward identities and Euclidean symmetry}
\label{ch14:ward}

We first use reflection positivity in each coordinate direction
separately. Continuous rotations are not assumed. The resulting
Hilbert spaces and bounded insertions provide the positivity estimates
needed to construct rotational Ward identities. Relativistic
reconstruction requires additional arguments and is not a consequence
of the Hilbert-space construction alone.

Throughout, $S$ is the mixed correlation hierarchy of
Chapter~\ref{ch13:continuum}. Its finite-dimensional field spaces are
$F_d$, with the compatible inclusions and derivative maps
constructed there. Statements using this calculus retain the actual
marked estimates required in that chapter. All correlations in the
present chapter are evaluated at distinct points or against test functions
supported away from collisions.

\section{Reflection spaces for the coordinate axes}
\label{ch14:reflection}

Fix an axis $\mu$. A positive word $w=((x_i,A_i))_{i=1}^n$ has
$x_i^\mu>0$ and distinct points; the empty word is allowed.
Let $\theta_\mu w$ reflect the points across $x^\mu=0$, apply the
signed coordinate reflection to the labels, conjugate them, and
reverse their order. On formal finite linear combinations set
\[
 \langle v,w\rangle_\mu=S(\theta_\mu v\,w).
\]
The form is positive semidefinite by reflection positivity. Its null
space is orthogonal to every vector by Cauchy--Schwarz. The quotient
and completion give $\mathcal H_\mu^{\rm OS}$, with dense word
vectors $[w]_\mu$ and unit vacuum $\Omega_\mu=[\varnothing]_\mu$.
This is the standard Osterwalder--Schrader quotient
\cite[Section3]{OsterwalderSchrader1973}. Only the quotient
construction is used here, not the reconstruction theorem.

Let $\tau_t$ mean translation by $te_\mu$. On formal word sums put
$T_tc=[\tau_tc]_\mu$, $t\ge0$. Translation invariance implies
\begin{equation}\label{ch14:heat-pairing}
 \langle T_sc,T_td\rangle_\mu
             =\langle T_0c,T_{s+t}d\rangle_\mu.
\end{equation}
Every fixed orbit $T_tc$ is bounded: the two reflected word halves
have a fixed positive separation and the separated-correlation
bound of Chapter~\ref{ch13:continuum} is uniform as they move apart.
If $M=\sup_{t\ge0}\|T_tc\|$, Cauchy--Schwarz in
\eqref{ch14:heat-pairing} gives $M^2\le\|T_0c\|M$.
Thus $\|T_tc\|\le\|T_0c\|$, including when $c$ is null.
The operators descend and extend to self-adjoint contractions
$E_\mu(t)$. Their semigroup law holds on the dense word span and
hence everywhere. Correlation continuity gives strong continuity.
The spectral theorem yields
\begin{equation}\label{ch14:heat}
 E_\mu(t)=e^{-tH_\mu},\qquad H_\mu\ge0,
 \qquad \lim_{t\to\infty}E_\mu(t)=P_\mu^0
 \quad\hbox{strongly},
\end{equation}
where $P_\mu^0$ projects onto $\ker H_\mu$.

Transverse translations, by $y$ with $y^\mu=0$, preserve the
reflected form and have inverse translations. They therefore give
strongly continuous unitaries $U_\mu(y)$, commuting with $E_\mu(t)$.
Both leave $\Omega_\mu$ fixed. At this stage we do not claim that
$\ker H_\mu$ is one-dimensional, or that its vectors are fixed by
transverse translations.

The commuting spectral measures of $H_\mu$ and the three transverse
momenta give a joint projection-valued measure $P_\mu$ on
$[0,\infty)\times\RR^3$:
\begin{equation}\label{ch14:coordinate-spectrum}
 E_\mu(t)U_\mu(y)
     =\int e^{-tE+ip\cdot y}\,dP_\mu(E,p).
\end{equation}
This is a coordinate half-space spectrum. No cone bound
$E\ge|p|$ has been proved or used.
For $\operatorname{Re}z>0$, the same spectral integral defines the
holomorphic contraction $e^{-zH_\mu}$; its derivatives are bounded
on every smaller right half-plane. This elementary heat analyticity
will be used before any rotation action is constructed.

\section{Bounded insertions in a slab}
\label{ch14:slabs}

A word $w$ lies in a slab of width $h$ if $0<x_i^\mu<h$ for each
of its points. Its reflected slab word is
$w^\sharp=\tau_h\theta_\mu w$. Formally define
$I_wv=w\tau_hv$. Reflection and translation give
\begin{equation}\label{ch14:formal-adjoint}
 \langle[I_wc]_\mu,[d]_\mu\rangle
       =\langle[c]_\mu,[I_{w^\sharp}d]_\mu\rangle.
\end{equation}
This identity is on formal word sums; neither boundedness nor
preservation of null vectors is assumed yet.

\begin{proposition}[Slab norm bound]\label{ch14:slab-bound}
For labels of bounded degree and norm, a word of length $n\ge1$,
with pairwise separation and slab margins at least $r$, $0<r\le1$,
defines a bounded operator $B_h(w)$ on $\mathcal H_\mu^{\rm OS}$.
There are constants depending only on the label bounds such that
\[
 B_h(w)[v]_\mu=[w\tau_hv]_\mu,
 \qquad \|B_h(w)\|\le C^n(n/r)^{\kappa n}.
\]
The constants are independent of the width and the exterior word.
\end{proposition}

\begin{proof}
Proposition~\ref{ch13:word-bound} first gives the estimate at a
width $s\le h$, with a slightly smaller buffer. Put
$T=I_{w^\sharp}I_w$ at that width and
$D=C^n(n/r)^{\kappa n}$. For every formal finite sum $c$ its
repeated-word estimate gives
$\|[T^jc]_\mu\|\le C_cD^{2j}$. By
\eqref{ch14:formal-adjoint}, $T$ is symmetric for the possibly
degenerate reflected form. Set
$b_j=D^{-2j}\|[T^jc]_\mu\|$ and $M_c=\sup_{j\ge0}b_j$.
Then $b_j^2\le\|[c]_\mu\|b_{2j}$, so
$M_c\le\|[c]_\mu\|$. In particular
\[
 \|[I_wc]_\mu\|^2
 =\langle[c]_\mu,[Tc]_\mu\rangle
 \le D^2\|[c]_\mu\|^2.
\]
The operator preserves null vectors and extends boundedly. Increase
the width to $h$ by right multiplication by $E_\mu(h-s)$.
\end{proof}

The empty word gives $B_h(\varnothing)=E_\mu(h)$. Dense word
matrix elements also prove the exact identities
\begin{align}\label{ch14:slab-laws}
 B_h(w)^*&=B_h(w^\sharp),\\
 B_{h+k}(w\tau_hv)&=B_h(w)B_k(v),\notag\\
 B_{h+t}(w)&=B_h(w)E_\mu(t),\qquad
 B_{t+h}(\tau_tw)=E_\mu(t)B_h(w).\notag
\end{align}
Transverse covariance reads
$U_\mu(y)B_h(w)U_\mu(y)^*=B_h(w+y)$.
These assertions concern bounded slab insertions, not unsmeared
point fields acting on an arbitrary Hilbert vector.

\begin{proposition}[Norm-smooth insertion positions]
\label{ch14:slab-smooth}
On the open configuration space inside a fixed slab, $B_h(w)$ is
smooth in operator norm. A position derivative replaces the
corresponding field label by its derivative in $F_{d+1}$.
\end{proposition}

\begin{proof}
For a fixed configuration choose a smaller common buffer. The norm
bound holds locally for the original tuple and each differentiated
tuple. On dense word matrix elements the derivative identity is
exactly the scalar derivative identity of Chapter~\ref{ch13:continuum}.
If a weak derivative $D A(x)[v]=C(x)[v]$ on a dense domain is locally
bounded in operator norm, the scalar mean-value inequality gives
$\|A(y)-A(x)\|\le\sup_{[x,y]}\|C\|\,|y-x|$.
Apply this first to every differentiated tuple, proving continuity
of its operator. The same inequality after subtracting
$C(x)[y-x]$ proves norm differentiability of $A$. Iteration proves
all orders. This argument explains why weak differentiability alone
would have been insufficient.
\end{proof}

\section{Comparing the mixed fields with free fields}
\label{ch14:comparison}

The comparison is a statement about the actual marked construction,
stronger than equality of dimensions of field spaces. We record its
quantitative form and then derive it from the finite-cutoff comparison. For fixed arity and degree bounds,
let $M_k\to\infty$, $r_k\to0$ and $\epsilon_k\to0$. Suppose the
same coefficient lifts in every slot give rescaled kernels $S_k$,
with, on a compact separated window,
\begin{equation}\label{ch14:grid-input}
 |S_k|\le A_k^n,\quad
 |\nabla S_k|\le C_n V_k^n/\delta_k,
 \quad |S_k-S^0|\le\epsilon_k\quad\hbox{on its grid}.
\end{equation}
The gradient bound is required outside coordinate-equality strips
of width $4\delta_k$; the grid has spacing $s_k$, and $S^0$ is the
free correlation. The necessary small errors are
\begin{equation}\label{ch14:grid-cost}
 A_k^n\delta_k\to0,\qquad V_k^n s_k/\delta_k\to0,
 \qquad s_k=o(\delta_k).
\end{equation}
For example the source construction uses
$r_k=a_kM_k^{\alpha_0}$, $s_k=M_k^{-\alpha_0}$,
$A_k=CM_k^{(\alpha_0-\alpha_1)D}$,
$V_k=CM_k^{(\alpha_0-\alpha_2)D}$ and $\delta_k=M_k^c$,
with strictly negative exponents in \eqref{ch14:grid-cost}.
The marked coefficient inverse, Gaussian comparison and displacement
estimates must all concern these same lifts. The following construction verifies this compatibility.

% !TEX root = ../monograph.tex
\subsection{Obtaining the three comparison bounds from one lattice family}
\label{ch14:lattice-comparison}

Here are the scale choices behind \eqref{ch14:grid-input}. Fix the
upper bound $n$ on the number of insertions, maximum scaling dimension $D$ and a compact set of configurations
$K$ with positive point separation. Enlarge the finite label list to
include its signed coordinate images and conjugates. All constants and
all three exponents below are chosen for this one finite request.
Choose positive numbers $\alpha_0,\alpha_1,\alpha_2$ and $c<0$ with
\begin{align}
 \alpha_1+6&<\alpha_0,&\alpha_2+5&<\alpha_0+c,\notag\\
 nD(\alpha_0-\alpha_1)+c&<0,&
 nD(\alpha_0-\alpha_2)-\alpha_0-c&<0.
 \label{ch14:comparison-exponents}
\end{align}
For example first choose $\alpha_0-\alpha_1=7$, next $-c>7nD$,
then $\alpha_0-\alpha_2>5-c$, and finally increase all three
$\alpha_i$ together until the last inequality holds and all are
positive. Increasing their common part does not alter the preceding
inequalities.

At a comparison scale $k$ whose comparison window lies in a single stage put\mgubi{Issue 10: say for which scales this holds ($k=k_j$ suffices for the limit taken below) and cite Chapter~\ref{ch10:trajectory}.}
$H_k=\exp((\log\log\widehat\beta_k)^2)$ and set
\[
 r_k=a_kH_k^{\alpha_0},\qquad s_k=H_k^{-\alpha_0},\qquad
 \delta_k=H_k^c.
\]
Thus $r_ks_k=a_k$: grid rounding occurs on the construction
lattice. Use one common localization map $T$ from
Section~\ref{ch11:common-extraction}. If $v$ is a fixed real polynomial
representative of a free quotient field, its lift $u_{k,N,p}$ solves
\[
                 S_{g/\theta}T u_{k,N,p}=W_{H_k^{\alpha_0}}v.
\]
Here $S_t$ multiplies a monomial with $f$ curvature factors by $t^f$,
whereas $W_\rho$ multiplies one of scaling dimension $j$ by $\rho^j$;
$g=\beta_{N,p,k}^{-1/2}$ and $\theta$ is the local normalization
used for $T$. These two dilations have different purposes.
The inverse bound in \eqref{ch11:simultaneous-detector-lifts} gives
bounded lifts at each fixed $k$, uniformly in the cutoffs and periods.
Their identities for reality, the identity field and signed symmetries are exact. No separate
choice is made for different spectator words.

For clarity, the use for finitely many insertions of the estimate for the corrected coefficients is
as follows. Write $I_i$ for the lattice insertion in argument $i$. For a subset
$J$ of the insertion indices, let $G_J$ be the free correlation of the
corresponding fields at the rescaled positions. Define recursively
\[
 R_\varnothing=1,\qquad
 R_I=\prod_{i\in I}I_i-
          \sum_{\varnothing\ne J\subseteq I}G_JR_{I\setminus J}.
\]
These are the Wick products of the insertions with respect to the free field (the moments $G_J$ are those of the free field, not of the lattice measure).
Equivalently, introduce commuting variables $s_i$ with $s_i^2=0$ and write
$s_J=\prod_{i\in J}s_i$. Then $R_I$ is the coefficient of $s_I$ in
$\prod_i(1+s_i I_i)/(1+\sum_{J\ne\varnothing}G_Js_J)$. The denominator has
constant coefficient one and hence an exact inverse in the finite
squarefree algebra. For the limiting Gaussian variables every
nonempty coefficient has expectation zero. The extraction of coefficients of source monomials,
covariance comparison and absorption of the larger source weight of
Proposition~\ref{ch11:n-mark-comparison} apply to this \emph{entire}
finite algebra, including its joint incoming terms.
With $q=\log H_k/(16n)$ and source weight $e^{2q}$, the output
normalization is at most $e^q$ per reflected cell, and extraction gives
\begin{equation}
 \|R_I\|_{\mathrm{cell}}
       \le e^{q-2q|I|}\le H_k^{-1/(16n)}
       \quad(\varnothing\ne I\subseteq\{1,\ldots,n\}).
 \label{ch14:subset-root}
\end{equation}
The lower bound on the reference coupling at the stage is chosen after this fixed number of insertions and the constants of its
insertions. The estimate is on the original positive partition
function and includes every reflected copy. Uniformity on $K$ uses
positive point separation in the covariance estimate; it does not
require a positive lower bound on each coordinate difference.
Unwinding the finite recursion and applying the chessboard inequality
therefore gives the grid error $C_KH_k^{-1/(16n)}$ at generic grid
points. Distinct coordinates there differ by at least $a_k$ in continuum
units, so the assignments to lattice sites require no rounding.

Apply the same coefficient map and the same finite comparison request
also at exponents $\alpha_1$ and $\alpha_2$. At its own normalization
a single insertion has chessboard seminorm at most $1+|G_{\{i\}}|$. Changing its
scaling dilation to $\alpha_0$ loses at most a factor
$CH_k^{(\alpha_0-\alpha_j)D}$. Set
\[
 A_k=CH_k^{(\alpha_0-\alpha_1)D},\qquad
 V_k=CH_k^{(\alpha_0-\alpha_2)D}.
\]
The comparison cells have sides at most $a_kH_k^{\alpha_j+5}$.
Consequently the first cells are $o(r_k)$, and the second are
$o(r_k\delta_k)$, by \eqref{ch14:comparison-exponents}.
The chessboard inequality gives $|S_k|\le A_k^n$ on $K$.
The modulus of continuity in the insertion points of Section~\ref{ch11:positional}, applied
with margin $r_k\delta_k$ in continuum units, gives
$|\nabla S_k|\le4nV_k^n/\delta_k$ outside the coincident-coordinate strips.
Include a strip for each configuration taking every coordinate from the point or
from its rounding (finitely many); the total volume of the strips is still $O(\delta_k)$.
The compact configuration sets in continuum units lie inside the smallest period for large $k$,
so the periodic and ordinary separations agree.

Take an ordered limit point of the lifts at each fixed $k$. The
uniform bounds pass to the already chosen continuum hierarchy.
Their exponents give exactly \eqref{ch14:grid-cost}, since
\[
 A_k^n\delta_k=CH_k^{nD(\alpha_0-\alpha_1)+c},\qquad
 V_k^ns_k/\delta_k=CH_k^{nD(\alpha_0-\alpha_2)-\alpha_0-c}.
\]
Both tend to zero, and $s_k/\delta_k\to0$ follows from the second
inequality. A diagonal sequence over the finite requests gives a
single comparison sequence. The bounded coefficient limits do not
alter the chosen correlation hierarchy or its order of limits.


The coordinate-equality strips have volume $O(\delta_k)$ on a fixed
window. Their contribution to any test function of supremum norm at most one
is $O((A_k^n+1)\delta_k)$. Outside them, rounding to the grid and
\eqref{ch14:grid-input} cost
$O(V_k^ns_k/\delta_k)$ plus the free modulus of continuity.
It follows that
\begin{equation}\label{ch14:uniform-tests}
 \sup_{\|f\|_\infty\le1,\ \operatorname{supp}f\subset K}
 \left|\int f(x)(S_k(x)-S^0(x))\,dx\right|\longrightarrow0.
\end{equation}
This is local $L^1$ convergence; it includes coordinate ties which
are not collisions. Testing against each fixed $f$ separately would
be weaker than the displayed assertion.

The lifts can be made exactly compatible with inclusions. Write
$i_d:F_d^0\to F_{d+1}^0$ for the free inclusion,
$j_d$ for the physical inclusion and $L_d$ for a real,
hypercubic-equivariant left inverse of $i_d$. From provisional maps
$U_{k,d}$ define recursively
\[
 \widetilde U_{k,d+1}
 =j_d\widetilde U_{k,d}L_d+U_{k,d+1}(I-i_dL_d).
\]
Exact inclusion identities of correlations preserve
\eqref{ch14:uniform-tests} in every slot. The unit is fixed at degree
zero. Free finite detecting families from Chapter~\ref{ch03:maxwell}
make these maps eventually injective; equality of dimensions in
Chapter~\ref{ch13:continuum} makes them isomorphisms. For each finite
request the tail is chosen once, before the test function. A single
sequence for all requests is obtained only by explicitly diagonalizing
arity, degree, compact windows and errors; the numerical comparison
level must be chosen after those finite requests.

\begin{proposition}[Derivative transport and smooth comparison]
\label{ch14:smooth-comparison}
For the comparison maps just described,
\[
 D_{k,d,a}:=r_kU_{k,d+1}^{-1}\partial_aU_{k,d}
                   \longrightarrow\partial_a^0
\]
in finite-dimensional operator norm. Rescaled correlations converge
in $C^j$ on every compact separated window, for every fixed $j$,
also for convergent moving labels.
\end{proposition}

\begin{proof}
Choose free tested correlations giving an injective map
$L:F_{d+1}^0\to\CC^q$. Their corresponding maps $L_k$
converge to $L$ by \eqref{ch14:uniform-tests}; hence
$\|L_kv\|\ge c\|v\|$ eventually. Integration by parts in the
insertion position gives
$L_k D_{k,d,a}v\to L\partial_a^0v$. No bound on
$D_{k,d,a}$ was assumed: subtracting the limiting derivative and
using the lower bound supplies it and its convergence.
Moving labels are handled by expansion in fixed finite bases.
Every fixed derivative of the correlation error consequently tends
to zero locally in $L^1$. Multiply by a smooth cutoff supported
away from collisions. Integration by parts gives a Fourier bound
by $(1+|p|)^{-M}$ times the sum of its $L^1$ derivative norms
through order $M$. Choose $M>4n+j$ and apply Fourier inversion.
\end{proof}

\section{The conserved stress line}
\label{ch14:stress-line}

Let $\mathcal K^0$ and $\mathcal K$ be the real spaces of centered,
real, symmetric, hypercubic-covariant degree-four tensor arrays in
the free and mixed field spaces. Centered means that each component
has zero expectation. Let $\mathcal V^0,\mathcal V$ be the degree-five
covariant vector spaces. The free divergence calculation of Chapter~\ref{ch03:maxwell}
has rank two. To check that its target is the full covariant vector
space, note that degree five permits only pairings of $F$ with $\partial F$ in $\gLie$. Sign parity and permutations leave the two
expressions
\[
 V_b^{(1)}=\sum_c{:}\langle F_{bc},\partial_bF_{bc}\rangle_{\gLie}{:},\qquad
 V_b^{(2)}=\sum_{c<d,\ c,d\ne b}{:}\langle F_{cd},\partial_bF_{cd}\rangle_{\gLie}{:}.
\]
The apparent further contractions
$\sum_{c,d}{:}\langle F_{bd},\partial_cF_{cd}\rangle_{\gLie}{:}$ and
$\sum_{c,d}{:}\langle F_{cd},\partial_cF_{bd}\rangle_{\gLie}{:}$ reduce to these: use the
Maxwell equation in the first and the Bianchi identity in the second.
The two independent rows of correlation pairings in \eqref{ch03:divergence-matrix}
show that these vectors are independent and in the divergence image.
Thus free divergence is onto, with kernel spanned by $T^0$.

Center the componentwise map $U_{k,4}$ after transport. It gives an
isomorphism $C_k:\mathcal K^0\to\mathcal K$: its inverse transports
back and centers there, and preservation of the unit makes the two
scalar corrections cancel. Let $V_k$ be the corresponding vector
isomorphism. The derivative comparison gives
\[
       V_k^{-1}r_k\operatorname{div}C_k
                          \longrightarrow\operatorname{div}^0.
\]
Compose with a fixed right inverse of free divergence. The result
tends to the identity and is invertible for large $k$, so physical
divergence is onto. Rank--nullity therefore gives
\begin{equation}\label{ch14:conserved-line}
                     \dim\ker(\operatorname{div}|_{\mathcal K})=1.
\end{equation}
This step uses the full low-degree rank computation, not just the
existence of a free conserved tensor.

Choose a nonzero $J$ on this line. Normalize its pulled-back arrays
$r_k^4 U_{k,4}^{-1}J$ to have the norm of $T^0$ and take a convergent
subsequence on that finite-dimensional sphere. Smooth comparison
at arity one ensures that its limit is centered. Derivative transport
ensures conservation. It must therefore be $T^0$ or $-T^0$.
Changing the sign of $J$ once gives positive numbers $c_k$ with
\begin{equation}\label{ch14:balanced-stress}
            U_{k,4}^{-1}(c_k^{-1}r_k^4J)\longrightarrow T^0.
\end{equation}
This fixes an orientation of the physical line. It does not yet fix
its Ward normalization, and no positive lower bound for $c_k$ has
been asserted.

\section{Shrinking shells and local charge fields}
\label{ch14:shells}

Choose $\chi\in C_c^\infty(\RR^4)$ equal to one near an insertion
$x_0$ and zero near all spectators. For a Euclidean Killing field
$v(y)$, either constant or $K(y-x_0)$ with $K^{\mathsf T}=-K$, set
\begin{equation}\label{ch14:flux}
 \mathcal Q^J_v(A;\boldsymbol B)
 =-\int\partial_\mu\chi(y)v_\nu(y)
              S(J_{\mu\nu}(y),A(x_0),\boldsymbol B)\,dy.
\end{equation}
Repeated indices are summed. Conservation and symmetry of $J$ show
that the value is independent of the cutoff as long as it separates
the same distinguished insertion from the spectators. Indeed the
difference is the integral of a total divergence in a collision-free
region, and $(\partial_\mu v_\nu)J_{\mu\nu}=0$. In particular one
can shrink a radial cutoff as in Figure~\ref{ch14:fig-ward-shell}.

\begin{figure}[htbp]
\centering
\begin{tikzpicture}[font=\small,
    inner shell/.style={draw=ymInk,semithick},
    outer shell/.style={draw=ymAccent,semithick,dashed}]
  % A central section of two radial cutoffs in four dimensions.
  \fill[ymAccent!13,even odd rule] (0,0) circle (1.80) (0,0) circle (1.20);
  \fill[ymInk!16,even odd rule] (0,0) circle (.99) (0,0) circle (.66);
  \foreach \r in {.66,.99} \draw[inner shell] (0,0) circle (\r);
  \foreach \r in {1.20,1.80} \draw[outer shell] (0,0) circle (\r);
  \fill (0,0) circle (1.6pt) node[below=2pt] {$A(x_0)$};
  \fill (-2.26,.76) circle (1.6pt) node[above] {$B_1(x_1)$};
  \fill (2.23,-.68) circle (1.6pt) node[below] {$B_2(x_2)$};

  % The horizontal scale here is twice that of the central section.
  \begin{scope}[shift={(3.60,-.85)}]
    \fill[ymInk!12] (1.32,0) rectangle (1.98,1.70);
    \fill[ymAccent!10] (2.40,0) rectangle (3.60,1.70);
    \draw[-{Stealth[length=1.7mm]}] (0,0) -- (5.02,0)
        node[above left=1pt] {$|y-x_0|$};
    \draw[-{Stealth[length=1.7mm]}] (0,0) -- (0,2.12);
    \node[left=2pt] at (0,1.70) {$1$};
    \node[left=2pt] at (0,0) {$0$};
    \draw[inner shell] (0,1.70) -- (1.32,1.70)
        .. controls (1.60,1.70) and (1.70,0) .. (1.98,0) -- (4.42,0);
    \draw[outer shell] (0,1.70) -- (2.40,1.70)
        .. controls (2.90,1.70) and (3.10,0) .. (3.60,0) -- (4.42,0);
    \node[anchor=east] at (1.53,.88) {$\chi_\rho$};
    \node[anchor=west] at (3.18,.88) {$\chi_R$};
    \foreach \r/\name in {1.98/\rho,3.60/R}
      \draw (\r,0) -- +(0,-.09) node[below=1pt] {$\name$};
  \end{scope}
\end{tikzpicture}
\caption[Deformation of a local Ward shell]{Shrinking the Ward shell. The two shaded annuli are planar
sections of the supports of $\nabla\chi_\rho$ and $\nabla\chi_R$.
The radial profiles agree near $x_0$ and near both spectators. Thus
$\chi_R-\chi_\rho$ is supported away from all three insertions, and
conservation gives the same flux on the two shells.}
\label{ch14:fig-ward-shell}
\end{figure}


The actual two-mark estimate needed here has the following form.
If $A\in F_d$, let $\varkappa=0$ for constant $v$ and
$\varkappa=1$ for linear $v$. Put $D=d+1-\varkappa$.
For a finite family of separated spectator correlation functionals $\ell_i$ on
$F_D$ and coefficients $g_i$ annihilating that space,
there must be $\rho_m\downarrow0$ such that
\begin{equation}\label{ch14:annular-input}
 \sup_{1/4\le|z|\le2}
 \left|\sum_i g_i\rho_m^{3+\varkappa}
       S(J_{\mu\nu}(x_0+\rho_m z),A(x_0),\boldsymbol B_i)\right|
                  \longrightarrow0.
\end{equation}
Here the same correction coefficient at the distinguished insertion
must work for every spectator in the finite family.
The power $3+\varkappa$ is the four-dimensional volume minus one
cutoff derivative plus the degree of the Killing field.

The mixed local replacement estimate of
Proposition~\ref{ch18:mixed-bound} supplies this assertion directly,
without using phase decomposition or rotational Ward identities.
Apply it to the two-field cluster $(J_{\mu\nu},A)$, of total degree
$d+4$, and truncate at $D=d+1-\varkappa$. Its replacement field
$C_D(\rho,z)\in F_D$ is chosen by fixed probes before the exterior
spectator. For every spectator in the finite family the residual is
bounded uniformly on $1/4\le|z|\le2$ by
$C\rho^{D+1-(d+4)-\epsilon}$. The same coefficients therefore cancel
under $\sum_i g_i\ell_i=0$. After multiplying by the shell factor,
\[
 \rho^{3+\varkappa}
 \left|\sum_i g_iS(J_{\mu\nu}(x_0+\rho z),A(x_0),\boldsymbol B_i)\right|
 \le C\rho^{1-\epsilon}\longrightarrow0,
 \qquad 0<\epsilon<1.
\]
This proves \eqref{ch14:annular-input} for every sufficiently small
radius, and hence along any sequence $\rho_m\downarrow0$.
The full compact annulus is allowed because the local replacement
proof already removed its lattice-grid restriction by derivative
bounds. Only point collisions are excluded; coordinate ties cause
no additional restriction. This use of local replacement precedes
its phase-uniform operator-product formulation.

\begin{proposition}[Representation of the local flux]
\label{ch14:charges}
There are unique linear maps
$p_b:F_d\to F_{d+1}$ and
$Y_d(K):F_d\to F_d$ such that the constant and
linear fluxes \eqref{ch14:flux} equal the correlations with
$p_bA$ and $Y_d(K)A$ at $x_0$, respectively.
\end{proposition}

\begin{proof}
Every finite linear dependence among the spectator functionals on
$F_D$ annihilates the flux: shrink the shell, rescale its
integration variable and use \eqref{ch14:annular-input}. Choose a
finite separating subfamily of spectator functionals. Its injective
map from $F_D$ has finite-dimensional range; the preceding
annihilator property says that the corresponding flux vector lies
in that range. It therefore specifies a unique field, and adjoining
any other spectator correlation functional proves the identity for it as well.
Continuity extends the identity from generic grid configurations to
every separated configuration. Linearity follows from uniqueness.
\end{proof}

Differentiating the shell identity with the distinguished insertion,
using conservation to move the cutoff and symmetry to remove the
Killing derivative, gives
\begin{equation}\label{ch14:charge-commute}
 Y_{d+1}(K)\partial_bA
   =\partial_bY_d(K)A+\sum_cK_{cb}p_cA,
 \qquad \partial_bp_cA=p_c\partial_bA.
\end{equation}
Equivalently one can differentiate after moving the whole local
shell with $x_0$; the derivative of $K(y-x_0)$ supplies the displayed
constant flux term. Faithfulness of separated spectator correlation functionals turns
these response identities into field identities. The same argument
shows compatibility with inclusions, reality and signed coordinate
symmetries; both charges annihilate the unit.

\section{Coordinate spectral measures and the angular boundary}
\label{ch14:spectral}

We recall that $P_\mu$ in \eqref{ch14:coordinate-spectrum} is the
joint measure of nonnegative heat energy and transverse momentum.
For a point field, translate its insertion by $ue_\mu$, $u>0$,
into the positive half-space and define the undamped measure
\[
 d\nu_\mu^A(E,p)=e^{2uE}
       d\langle[A(ue_\mu)],P_\mu(E,p)[A(ue_\mu)]\rangle.
\]
The semigroup law makes this independent of $u$. It is finite on
$E\le M$; it need not have finite total mass. Word measures are
defined after translating the entire word beyond its extreme point.
Polarization defines mixed measures. For integrable nonnegative
$g_1,g_2$ and $|f|^2\le g_1g_2$, spectral Cauchy--Schwarz gives
\begin{equation}\label{ch14:spectral-CS}
 \left|\int f\,d\nu_\mu^{AB}\right|^2
       \le\int g_1\,d\nu_\mu^A\int g_2\,d\nu_\mu^B.
\end{equation}
Positive damping and the separated derivative bounds imply all
polynomial transverse moments. One direct justification takes
centered differences of the reflected two-point function: their
spectral multiplier is $(2\sin(tp_j/2)/t)^{2r}$. Fatou's lemma
and the integral remainder for finite differences give the moment
bound as $t\to0$. Further heat damping gives energy moments.
For point fields, integration over small damping times against a
sufficiently high power of time also gives
\begin{equation}\label{ch14:polynomial-moments}
       \int (1+E)^{-D}|p_j|^{2r}\,d\nu_\mu^A<\infty
\end{equation}
for some $D$ depending on $A,r$.

\subsection{What opposite reflection signs control}

If a field or word combination $A$ has reflection signs
$\varepsilon_i=\pm1$ in axes in a set $I$, its two-point kernel
has, in time direction $i$, the spectral sign $\varepsilon_i$.
For an even frequency profile $f$ use
\[
 \mathcal Lf(E)=\int e^{-E|t|}\left(\int e^{ipt}f(p)\,dp\right)dt
       =2E\int\frac{f(p)}{p^2+E^2}\,dp.
\]
At $E=0$ the limiting value is $2\pi f(0)$.
To compare representations in different axes without an arbitrary
collision extension, subtract the low moments of the profiles.
Choose even $\sigma\ge0$ of integral one supported in
$1\le|v|\le2$, and the Lagrange polynomials $L_j$ for the nodes
$(j+2)^2$, $0\le j<m$. Define the annular correction by
\[
 \int (\mathcal C_{m,\rho}f)(q)g(q)\,dq
 =\int\sigma(v)f(p)\sum_{j<m}L_j((p/(\rho v))^2)
                           g(\rho v(j+2))\,dp\,dv.
\]
Then $f-\mathcal C_{m,\rho}f$ has its first $2m$ moments zero.
Interpolation of $(p^2+E^2)^{-1}$ gives its exact response
\begin{equation}\label{ch14:positive-residual}
 \mathcal R_{m,\rho}f(E)
 =2E\int\frac{f(p)}{p^2+E^2}\int\sigma(v)
       \prod_{j<m}\frac{(\rho v(j+2))^2-p^2}
                          {(\rho v(j+2))^2+E^2}\,dv\,dp.
\end{equation}
For $f\ge0$ supported in $[-\rho,\rho]$ this is nonnegative.
If $f_w(p)=w^{-1}f(p/w)$ with $f=1$ near zero and
$0<w\le1\le\rho$, it is bounded below by
$c_m/(E+w)$ on $0\le E\le1$. For large $E$ its decay can be made
arbitrarily high by increasing $m$. The upper bound, unlike the
low-energy lower bound, holds also for growing $w$: it requires only
$\rho\ge1$ and $\operatorname{supp}f_w\subset[-\rho,\rho]$.
For $E\le1$ use $\|f_w\|_\infty\le C/w$; for $E\ge1$ the
product in \eqref{ch14:positive-residual} is bounded by
$C_m\min(1,(\rho/E)^{2m})$. Thus
$\mathcal R_{m,\rho}f_w(E)\le C(w^{-1}+\rho^D)(1+E)^{-D}$
for any fixed $D\le2m$. Together with
\eqref{ch14:polynomial-moments} this gives, for fixed $r$,
\[
 \int\mathcal R_{m,\rho}f_w(E)|p_j|^{2r}\,d\nu_i^A
                                      \le C(w^{-1}+\rho^D).
\]
Choose $D$ from the damped moment first, and then $m$ sufficiently
large; $D$ does not grow with the later subtraction order.

We now give the precise comparison. For a word combination $A$, assume
its offsets have zero coordinates in $I$ and that reflection in axis $i$
acts by $\varepsilon_i$ for $i\in I$. These geometric restrictions permit
arbitrarily small positive placements in any of the chosen time axes,
so the integrated moment bound \eqref{ch14:polynomial-moments} holds
in each. They will be relaxed for the line axis below.
For even nonnegative fixed profiles $f_j\in C_c^\infty(\RR)$ write
$f_{j,v}(p)=v^{-1}f_j(p/v)$; in particular the normalization preserves
its integral, while its supremum costs $v^{-1}$. For a product of profiles
$g=(g_k)_{k\in I}$ define
\[
 \operatorname{Resp}_i(g)
   =\int\mathcal Lg_i(E)\prod_{k\in I\setminus\{i\}}g_k(p_k)
                                                     \,d\nu_i^A(E,p).
\]
The same even product test function of the unreflected pair kernel has spectral
values $\varepsilon_i\operatorname{Resp}_i(g)$ in every axis $i$.
First require sufficiently many vanishing moments in two axes, so the
test function vanishes to the required order at the collision. This makes the
comparison an equality of convergent integrals of the original kernel.
If $I$ omits axes, insert in each omitted frequency coordinate a profile
$h(p/T)$ with $h=1$ near zero, and let $T\to\infty$; the damped moment
bounds justify this limit. Thus opposite signs give
$\operatorname{Resp}_i(g)=-\operatorname{Resp}_j(g)$.

Let $a_i=\mathcal C_{m,\rho_i}g_i$ and
$a_j=\mathcal C_{m,\rho_j}g_j$, where each radius contains the support
of its profile and $m$ exceeds the necessary collision orders.
Apply the signed equality to the product with residual profiles in
both $i$ and $j$. Expanding it gives
\begin{align}
 \operatorname{Resp}_i(g_i-a_i;g_{\ne i})
 &\le \operatorname{Resp}_i(g_i-a_i;g_j\mapsto a_j)
       +\operatorname{Resp}_j(g_j-a_j;g_i\mapsto a_i).
 \label{ch14:sign-transfer}
\end{align}
Here the semicolon leaves every unmentioned profile unchanged.
The omitted term is
$-\operatorname{Resp}_j(g_j-a_j;g_{\ne j})\le0$, because its temporal
residual response and all its transverse profiles are nonnegative.
The corrections $a_i,a_j$ themselves can have signs. They occur only
as transverse profiles on the right, so replacing them by nonnegative
majorants gives an upper bound. This explains precisely where positivity
is used.

We claim that when both signs occur in $I$, for $\lambda\in e^{\NN}$,
$q<0$, $\beta_j\le0$, and every $\delta>0$,
\begin{equation}\label{ch14:mixed-sign}
 \int_{E\le1}\frac{\prod_{j\in I\setminus\{i\}}
                         f_{j,\lambda^{\beta_j}}(p_j)}
                        {E+\lambda^q}\,d\nu_i^A
                         =O_\delta(\lambda^\delta).
\end{equation}
The following induction proves the claim while keeping the order of
parameter choices explicit. Choose a walk $i_0,\ldots,i_N$ through $I$
with opposite signs at consecutive entries, starting at $i_0=i$, with
every axis visited before its final entry. Put
$V_s=\{i_0,\ldots,i_{s-1}\}$. Fix $K\ge0$ bounding the magnitudes of
all negative width exponents. Induct backwards on $s$ to find $a_s>0$
such that
\begin{equation}
 \int\mathcal R_{m,R_0\lambda^\eta}
            f_{i_s,\lambda^{\omega_{i_s}}}(E)
          \prod_{j\ne i_s}f_{j,\lambda^{\omega_j}}(p_j)\,d\nu_{i_s}^A
                       \le C\lambda^\delta.
 \label{ch14:sign-induction}
\end{equation}
Here $0\le\omega_j\le a_s$ on $V_s$,
$-K\le\omega_j\le a_s$ elsewhere,
$0\le\eta\le a_s$ and $\omega_{i_s}\le\eta$.
The fixed $R_0\ge1$ contains the support at the current radius;
$m$ may be any sufficiently large fixed subtraction order.
Products here and below range over $I$.

At $s=N$ all widths are at least one, hence the transverse suprema
are bounded. The residual moment estimate with moment order zero
bounds the left side by $C(\lambda^{-\omega_{i_N}}+\lambda^{\eta D})$.
Choose $a_ND<\delta$. For the step from $s+1$ to $s$ write
$i=i_s$, $j=i_{s+1}$ and $p_*=K|I\setminus\{i,j\}|$.
Choose, in this order,
\[
 \alpha=a_{s+1},\qquad
 2r\alpha>p_*+\max(K,\delta)+1,\qquad
 D=D(A,r),\qquad
 0<a_s<\min\{\alpha,\delta/(D+1)\}.
\]
Only after these choices enlarge $m$ to handle the moment exponent $D$
and the collision orders. Correct $g_i$ at radius $R_0\lambda^\eta$
and $g_j$ at a fixed multiple of $\lambda^\alpha$.
The correction in $i$ has a fixed nonnegative annular majorant of
width $\lambda^\eta$, since its coefficients involve only factors
$\lambda^{2r'(\omega_i-\eta)}\le1$. The second term of
\eqref{ch14:sign-transfer} is therefore covered by the next induction
step, with the previously current axis now having nonnegative width
exponent. The correction in $j$ is bounded by
$C_r\lambda^{-2r\alpha}|p_j|^{2r}$ on its annulus; its coefficients
are bounded because $\omega_j<\alpha$. The remaining narrow suprema
cost at most $\lambda^{p_*}$. The first term is consequently at most
\[
 C\lambda^{p_*-2r\alpha+
                    \max(-\omega_i,\eta D)}
 \le C\lambda^{p_*-2r\alpha+\max(K,a_sD)}
                  =O(\lambda^\delta).
\]
This proves \eqref{ch14:sign-induction}. Finally take $s=0$,
$\omega_i=q$, $\omega_j=\beta_j$, and $\eta=0$.
The lower bound $\mathcal R_{m,R_0}f_{i,\lambda^q}(E)
\ge c_m/(E+\lambda^q)$ on $E\le1$, with a bump $f_i=1$ near zero,
gives \eqref{ch14:mixed-sign}. No continuous rotation was used.

\subsection{The separate diagonal-stress estimate}

For $J_{\mu b}$ with $\mu\ne b$, the two odd tensor indices give
opposite signs and \eqref{ch14:mixed-sign} applies. A diagonal
stress component has only positive signs, so it requires conservation.
We describe the contact subtraction needed just for this argument;
it is not a positive extension of the correlation hierarchy.

Let $K_{ab,cd}(x)=S(J_{ab}(x),J_{cd}(0))$ for $x\ne0$, and choose
an integer $D$ above the collision orders of these kernels and their
first derivatives. For one-dimensional Schwartz test functions let $\Pi_D$
retain Taylor orders less than $D$, multiplied by a fixed even
cutoff. On product test functions define
\[
 R_{ab,cd}(h)=\int K_{ab,cd}(x)
       \left(\prod_i h_i(x_i)-\prod_i\Pi_Dh_i(x_i)\right)dx.
\]
Each term of the telescoping product difference has a factor
$O(|x_i|^D)$, ensuring integrability. Using the same cutoff in every
axis preserves signed tensor symmetries and exchange of the index
pairs. Conservation makes its divergence a distribution supported at
zero with finitely many derivatives. Inversion removes its even orders.
Here is an explicit subtraction of the possible linear orders. Put
$\chi_0(x)=\prod_j\chi(x_j)$, using the even Taylor cutoff above, and set
\[
 q_{ab,cd}=\sum_i R_{ib,cd}\bigl(\partial_i(x_a\chi_0)\bigr),
 \qquad R'_{ab,cd}=R_{ab,cd}-q_{ab,cd}\delta_0.
\]
Indeed $\sum_i\partial_iR_{ib,cd}$ evaluated on $x_a\chi_0$ equals
$-q_{ab,cd}$, whereas the divergence of $q_{ib,cd}\delta_0$ has the
same value. Thus their difference has no linear derivative of
$\delta_0$. To check the tensor symmetries, reflection parity makes
$q_{ab,cd}$ vanish unless each coordinate index occurs an even number
of times. The remaining index patterns are $aa,aa$, $aa,bb$ and
$ab,ab$ (including their within-pair reversals); coordinate permutations,
within-pair symmetry of $R$, and exchange of its two pairs give the
same symmetries for $q$. The subtraction is therefore compatible with
all the tensor identities. We henceforth rename $R'$ as $R$.
Its Fourier divergence polynomial starts in degree three.

Write $\widehat R$ for this corrected functional and
$F_{b,cd}=\sum_i p_i\widehat R_{ib,cd}$. For a chosen time axis
$\mu$, set
\begin{align*}
 Q_0&=\widehat R_{\mu\mu,\mu\mu},\\
 Q_1&=\sum_{j,k\ne\mu}p_jp_k\widehat R_{k\mu,j\mu},\\
 Q_2&=\sum_{j,k,r,s\ne\mu}p_jp_kp_rp_s\widehat R_{sk,rj}.
\end{align*}
Conservation on the two index pairs gives
$Q_1=p_\mu^2Q_0+D_1$ and $Q_2=p_\mu^2Q_1+D_2$, where the
polynomials $D_1,D_2$ start in degrees four and six. Sufficiently
many temporal moment subtractions annihilate them exactly.
On a fixed temporal annulus the transverse responses of $Q_0,Q_1,Q_2$
are bounded by $O_\delta(\lambda^\delta)$ times respectively
$1,W^2,W^4$, if the transverse widths are at most $W$. For a diagonal
entry, divide its annular temporal profile by $p_\mu$ and put that
derivative on the label; the resulting vector has mixed signs.
Odd profiles are corrected by
$p^{-1}\mathcal C_{m,\rho}(pf)$, and their residual Laplace absolute
value is bounded by $C(E+w)^{-1}(1+E)^{-2m}$.
Thus \eqref{ch14:mixed-sign} and spectral Cauchy--Schwarz apply.
Removing one transverse moment restriction at a time terminates
after three axes. The annular coefficients remain bounded because
all shrinking width exponents are nonpositive.

For a fixed even annular $\sigma$, evaluate $Q_r$ against
$g_{v,r}(p)=v^{-1}(p/v)^{-2r}\sigma(p/v)$, so
$p^2g_{v,r+1}=v^2g_{v,r}$. Compare the widths $W$ and $2W$ and
subtract their fixed annular corrections. Write $c_r(t)=W^{-2r}\operatorname{Re}Q_r(g_{tW,r})$, with the
transverse profiles fixed, and let $d_r$ be the same normalized real
response to $(I-\mathcal C_{m,8})(g_{W,r}-g_{2W,r})$.
Conservation gives $c_{r+1}(t)=t^2c_r(t)+O_\delta(\lambda^\delta W^2)$
for $t=1,2$. The correction has zero constant moment and higher even
moments $O(W^2)$, so
$d_r=c_r(1)-c_r(2)+O_\delta(\lambda^\delta W^2)$. Eliminating
$c_0(1),c_0(2)$ from these three relations gives
\[
                  4d_0=5d_1-d_2+O_\delta(\lambda^\delta W^2).
\]
The corrected dyadic profile is signed, so its positivity needs
one further observation. For fixed correction nodes
$A_j=(8v(j+2))^2$, the function
\[
 R_E(z)=\frac1{z+E^2}\prod_{j<m}\frac{A_j-z}{A_j+E^2}
\]
is positive and decreasing for $0\le z\le16$, since $|v|\ge1$
and $A_j\ge256$. The dyadic response is the integral of
$2E[R_E(W^2p^2)-R_E(4W^2p^2)]$ against the nonnegative
annular profile $p^{-2r}\sigma(p)$ and the correction profile.
For $W\le1$ it is therefore nonnegative. On $E\le1$ the
product factors have a fixed positive lower bound; retaining the
difference of the two resolvents gives the lower bound
$cEW^2/(E+W)^4$ for $r=0$.
Conservation with the sufficiently many zero moments now gives
$d_0\ge0$, $d_1\le0$, $d_2\ge0$: each temporal multiplier
$p_\mu^2$ becomes $-E^2$ in the Laplace response, with the
contact terms killed exactly. Therefore, for
a nonnegative transverse product profile $\mathcal K_\lambda$,
\begin{equation}\label{ch14:diagonal-mass}
 \int_{E\le1}\frac{EW^4}{(E+W)^4}\mathcal K_\lambda(p)
                         \,d\nu_\mu^{J_{\mu\mu}}
                  =O_\delta(\lambda^\delta W^4).
\end{equation}
The integrand vanishes at $E=0$. The zero-energy part must be kept
separate in the subsequent Cauchy--Schwarz estimate.

\subsection{Vanishing angular flux at infinity}
\label{ch14:collinear}

For $K=E_{ba}-E_{ab}$ and spectator positions on a coordinate line
$\RR e_l$, choose anisotropic exponents
$\omega_a=\omega_b=-100$, $\omega_l=-150$ when $l\notin\{a,b\}$,
and $\omega_j=-300$ on remaining axes. Let
$w_j=\lambda^{1+\omega_j}$,
$\mathcal K_\lambda=\prod_{j\ne\mu}w_j^{-1}\psi(p_j/w_j)$,
with fixed even nonnegative $\psi=1$ near zero, and put
\[
 I_\mu(A;q)=\int_{E\le1}\frac{\mathcal K_\lambda(p)}
                                  {E+\lambda^q}\,d\nu_\mu^A.
\]
Project the spectator word to odd parity in $a,b$ and even parity
in the other axes, calling the result $B$. For a reflected line
axis this is a finite combination of words on the same line.
The mixed-sign estimate gives
\[
 I_\mu(B;-3000)=O(\lambda^{\gamma_\mu}),\quad
 I_\mu(J_{\mu b};\omega_\mu)=O_\delta(\lambda^\delta)\quad(\mu\ne b),
 \quad\gamma_l=1,\quad\gamma_\mu\le150.
\]
For the line time axis this requires a genuine temporal gap.
Choose a nonnegative real even smooth profile $\sigma_0$, let
$c=\sigma_0*\sigma_0$, and normalize
$\Phi=\widehat c=(\widehat\sigma_0)^2$ so $\Phi(0)=1$.
Thus $\Phi\ge0$. If translating the word by $u$ puts it in the
positive half-space, choose an even cutoff $\zeta=1$ on
$[-2u-1,2u+1]$, supported in $[-2u-2,2u+2]$.
Remove from $c_w$ the profile with Fourier transform
$\zeta(t)\Phi(wt)$. Its inverse-Fourier supremum is bounded
uniformly in $w$, whereas the residual has positive response
\[
 \int e^{-E|t|}(1-\zeta(t))\Phi(wt)\,dt
                          \ge c_u/(E+w)\quad(0\le E\le1).
\]
For the lower bound integrate on
$R<t<R+\gamma/(E+w)$, with $R=2u+2$, $wR\le\gamma$ and
$\Phi\ge1/2$ on $[-3\gamma,3\gamma]$.
Apply the sign-transfer inequality using this residual and an
opposite-sign transverse axis. Its correction term is bounded by
the three-axis mixed-sign estimate and the uniform supremum above;
its other term is bounded using a sufficiently high transverse moment,
now with fixed damping $e^{-2uE}$. The parity pattern has both signs
among the transverse axes and a sign opposite to the line axis,
so this proves any positive exponent, in particular $\gamma_l=1$.
For a different time axis use the three transverse-to-line axes
and bound the omitted line profile by its supremum. With width
$\lambda^{1+\omega_l}$ this costs exponent
$-(1+\omega_l)+1=-\omega_l\le150$.
These arguments keep the word measures fixed and include their
zero-energy mass; an arbitrary frequency bump would not have the
positive temporal-gap response used here. The diagonal estimate is
\eqref{ch14:diagonal-mass} with $W=\lambda^{-99}$.
The crude bound $I_\mu(J_{\mu\mu};0)=O(\lambda^{900})$ follows
from finite mass on $E\le1$ and the profile suprema. Outside
$E\le1$ and $|p_j|\le w_j$, positive damping and Schwartz decay
make the relevant test functions $O(\lambda^{-M})$ for every fixed $M$.

Take even $\eta\in C_c^\infty(\RR)$ equal to one near zero,
put $\ell(t)=t\eta(t)$, and define
\[
 f_{ab}(y)=\lambda^{100}\ell(\lambda^{-100}y_a)
                \prod_{j\ne a}\ell'(\lambda^{\omega_j}y_j),
 \qquad V_b=f_{ab},\quad V_a=-f_{ba}.
\]
The other components of $V$ vanish. Near all insertions $V(y)=Ky$.
Conservation and the local shell identities express the total raw
angular response as $-\int\partial_\mu V_c S(J_{\mu c},\boldsymbol A)$.
The terms $\partial_af_{ab}J_{ab}$ and
$\partial_bf_{ba}J_{ba}$ cancel exactly. In the remaining terms,
reflection projects the spectators to $B$. Fourier scaling supplies
$\lambda^3\mathcal K_\lambda$ and the temporal derivative supplies
$e^{-E/\lambda^{\omega_\mu}}$.
Apply \eqref{ch14:spectral-CS}. For the diagonal term split at
$E=\lambda^{-3000}$, using \eqref{ch14:diagonal-mass} above this
value and the crude bound below it. The worst resulting exponents,
including the prefactor $\lambda^{100}$, are
\[
\begin{array}{c|c}
 \mu=l\notin\{a,b\}&100+3-150+(1+1)/2=-46\\
 \mu\notin\{a,b,l\}&100+3-300+(150+1)/2=-121.5\\
 \mu=b,\ E\ge\lambda^{-3000}&100+3+(150+1-396)/2=-19.5\\
 \mu=b,\ E<\lambda^{-3000}&100+3-1500+(150+900)/2=-872.
\end{array}
\]
The exchanged column has the same bounds. Every term tends to zero,
whereas the total shell response is independent of $\lambda$.
Consequently
\begin{equation}\label{ch14:line-ward}
 \sum_i\left\{S(\ldots,Y(K)A_i,\ldots)
       +(Kx_i)_bS(\ldots,p_bA_i,\ldots)\right\}=0
\end{equation}
for every separated configuration on a coordinate line.
This argument used coordinate positivity and conservation, not a
cone spectrum or a pre-existing rotation action.

\section{Normalization and propagation of the Ward identities}
\label{ch14:propagation}

Translate a collinear configuration in its line direction in
\eqref{ch14:line-ward} and subtract the original identity. With a
plane generator mixing that line and a transverse direction one
obtains the translation Ward sum for $p_b$. The commutation law
\eqref{ch14:charge-commute} gives its positional derivatives as well.
The usual positional derivatives $\partial_b$ obey the same identities.

The finite vector correlation functionals of Theorem~\ref{ch03:vector-tests} have a
one-dimensional common kernel in the free vector-map space, spanned
by $\partial^0$. Adjoin a linear functional nonzero on that vector;
the matrix of correlation pairings becomes injective. Smooth comparison transports
this injectivity to the mixed finite-dimensional vector-map space.
Thus its Ward kernel has dimension at most one. It contains both
$p$ and $\partial$, so $p=c_d\partial$ at every sufficiently large
degree. Inclusions and nonzero free differentiation identify these
constants across degrees and extend the identity to the lower ones.
Call the common constant $c_*$. Changing variables in the shell
integrals and using \eqref{ch14:balanced-stress} gives
\begin{equation}\label{ch14:charge-limits}
 \frac{r_k}{c_k}U_{k,d+1}^{-1}p_bU_{k,d}\to\partial_b^0,
 \qquad
 \frac1{c_k}U_{k,d}^{-1}Y_d(K)U_{k,d}\to L_d^0(K).
\end{equation}
A finite separating family of correlation pairings justify these operator limits, even if the
operator norms were not known in advance. Since
$r_kU_{k,d+1}^{-1}\partial_bU_{k,d}\to\partial_b^0\ne0$
in some degree, the first limit gives $c_*\ne0$ and $c_k\to c_*$.
In particular $c_*>0$, since $c_k>0$.
Put $L_d(K)=c_*^{-1}Y_d(K)$. Then
\begin{equation}\label{ch14:derivative-commutator}
 L_{d+1}(K)\partial_b=\partial_bL_d(K)+\sum_cK_{cb}\partial_c.
\end{equation}
Every positional derivative of the normalized angular Ward sum is
therefore another such sum with a differentiated label.

\subsection{Analyticity supplied by a positive Gram kernel}

The missing step from a line to an open set is analyticity, not
smoothness. We establish exactly the analyticity needed.
For $x^a>0$ let $\Psi_A(x)=[A(x)]_a$ in the coordinate reflection
space. The curve $t\mapsto\Psi_A(x+te_b)$ is smooth for $b\ne a$:
its weak derivative against every word vector is
$\Psi_{\partial_b A}(x+te_b)$, which is continuous by translation
continuity. Integrating this continuous vector curve and using the
dense word span proves strong differentiability. Iterate for higher
orders. In the normal direction the curve is analytic away from the
reflection plane, since it is a positive-time heat orbit.

Let $F_{A,B}(x)=S(A(x),B(0))$. Its restriction to $ue_a$, $u>0$,
is analytic by that heat argument. In the plane $(a,b)$ with
$Ke_a=e_b$, $Ke_b=-e_a$, define
\[
 Q_{A,B}(u,z)=F_{e^{-\theta L(K)}A,e^{-\theta L(K)}B}(\rho e_a),
 \quad \rho=(u^2+z^2)^{1/2},\quad\theta=\arctan(z/u).
\]
This is analytic for $u>0$: the label spaces are finite-dimensional
and the radial functions are analytic. It solves the angular
identity because
$(u\partial_z-z\partial_u)\theta=1$ and the same operator kills
$\rho$. It agrees with the actual $F$ at $z=0$.
All transverse jets of the actual angular expression vanish there,
by \eqref{ch14:line-ward} and
\eqref{ch14:derivative-commutator}. For $H=F-Q$, differentiating
the angular expression $n$ times in $z$ gives at $z=0$
\[
 u\partial_z^{n+1}H_{A,B}
 =-\partial_z^n(H_{L(K)A,B}+H_{A,L(K)B})
                        +n\partial_u\partial_z^{n-1}H_{A,B}.
\]
The last term is omitted at $n=0$. Induction proves equality of all
transverse jets of $F$ and $Q$, simultaneously for all labels.

Equality of jets has not yet proved equality of functions. The
positive Gram identity supplies the decisive bound. If a smooth
Hilbert-valued curve $v(t)$ satisfies
$\langle v(s),v(t)\rangle=A(t-s)$ and the derivatives of $A$ at
zero agree with an analytic function, then
\[
 \|v^{(n)}(t)\|^2=|A^{(2n)}(0)|
 \le C R^{-2n}(2n)!
 \le C(2/R)^{2n}(n!)^2.
\]
Hilbert-space Taylor remainders tend to zero for $|h|<R/2$,
uniformly in $t$. Apply this to the single-insertion curve above: its Gram
kernel is $F_{A,A^{\dagger_a}}(2x^ae_a+(t-s)e_b)$, whose jets
are those of $Q$. Thus single-insertion vectors are analytic along every
transverse coordinate as well.

If one insertion is a strict maximum in some coordinate, separate
it from the spectator word by a coordinate hyperplane. Its moment
is a matrix element of this analytic single-insertion curve against a
fixed word vector. Strict minima are handled by the existing signed
reflection. Hence the Ward sum is analytic along any coordinate
segment on which the moving point stays a strict extremum in some
coordinate. This uses no integration of the spin operators into a
rotation group.

\subsection{A finite path through separated configurations}

Configurations with distinct first coordinates are dense. Starting
from one, choose the point with smallest first coordinate among
those not yet moved. Move it in the second coordinate until it
exceeds all other second coordinates; it remains strictly first in
the first coordinate. Next move it in the first coordinate beyond
all other points, remaining extreme in the second. Finally set its
other coordinates to zero, remaining extreme in the first. Repeat
with the remaining points. Every segment is separated and of the
analytic kind just proved. The result lies on the first coordinate
axis. There all Ward sums and all their derivatives vanish.
Reverse the path: the one-variable identity theorem on each segment
propagates the identity and its derivative identities. Continuity
then gives, on every separated configuration,
\begin{equation}\label{ch14:global-ward}
 \sum_i\left\{S(\ldots,L(K)A_i,\ldots;x)
                  +(Kx_i)_b\partial_{x_i^b}S(\boldsymbol A;x)\right\}=0.
\end{equation}
This is why replacing analyticity by smoothness in the propagation
argument would leave a gap.

\section{The rotation action and the normalized tensor}
\label{ch14:rotations-section}

Commuting two identities \eqref{ch14:global-ward} shows that the
internal defect
$[L(K),L(M)]-L([K,M])$ has zero correlation sums. It is real,
alternating, filtration preserving, kills the unit, and has the
signed covariance of a pair of rotation generators. The finite
internal-defect correlation functionals of Theorem~\ref{ch03:internal-tests} are
injective on this defect space. Their injectivity transfers under
the smooth finite-dimensional comparison, exactly as for the vector
correlation functionals above. The defect is therefore zero. Injective inclusions
extend this conclusion from any sufficiently large degree to the
lower ones.

One must still exclude a representation only of the universal cover
of $SO(4)$. Here the free comparison gives an exact intertwiner.
Choose the six plane generators and, at one degree, write their free
matrices as $C_i$ and their transported mixed matrices as $B_{k,i}$.
By \eqref{ch14:charge-limits} these converge to $C_i$. On the matrix
space put
\[
 \mathcal A_{k,i}(X)=B_{k,i}X-XC_i,
 \qquad \mathcal C_k=\sum_i\mathcal A_{k,i}^2.
\]
The bracket identities make $\mathcal C_k$ commute with every
$\mathcal A_{k,i}$. At the limit an invariant inner product makes
$\mathcal C$ self-adjoint and negative semidefinite, with kernel
the commutant of the free action. Its zero eigenvalue is isolated.
Finite-dimensional spectral projections around that eigenvalue give
real projections $P_k\to P$ commuting with all $\mathcal A_{k,i}$.
They also commute with the signed symmetries on matrix space.
The restrictions of $\mathcal A_{k,i}$ to $\operatorname{ran}P_k$
tend to zero. Every plane generator is a commutator of two others,
so their maximal norm satisfies $m_k\le2m_k^2$. Eventually $m_k=0$.
Consequently $f_k=P_kI\to I$ is invertible and
$B_{k,i}f_k=f_kC_i$. The intertwiner $U_{k,d}f_k$ is real and
preserves the signed action. Conjugating the free $O(4)$ action by
it gives the required genuine orthogonal-group representation.
This finite-dimensional argument also fixes the periodicity of
plane rotations, rather than assuming it from Lie brackets.

\begin{theorem}[Euclidean rotations]\label{ch14:rotations}
Under the actual marked comparison and annular estimates used above,
there are continuous complex-linear representations $R_d$ of $O(4)$
on $F_d$, extending the signed coordinate symmetries,
preserving reality and the unit, and compatible with inclusions.
Their infinitesimal generators are $L_d$, their derivatives transform
as vectors, and
\[
 S(R_{d_i}(g)A_i;gx)=S(\boldsymbol A;x),\qquad g\in O(4).
\]
They are isomorphic in each degree to the free label representations.
\end{theorem}

\begin{proof}
The preceding construction supplies the representations. The function
\[
 t\longmapsto S(R(e^{tK})\boldsymbol A;e^{tK}x)
\]
has zero derivative by \eqref{ch14:global-ward}. Plane rotations and
signed reflections
generate $O(4)$. Covariance, the fact that inserting the identity leaves correlations unchanged,
and the exact derivative and inclusion identities show that the corresponding label identities
have zero responses against every spectator. Faithfulness of the
field quotient proves their equality. Thus the independently chosen
finite-degree intertwiners need not themselves commute with inclusions;
the resulting actions do.
\end{proof}

Let $J_0=c_*^{-1}J$. With probability Haar measure on $O(4)$ set
\begin{equation}\label{ch14:averaged-stress}
 T_{ab}=\int_{O(4)}\sum_{i,j}g_{ai}g_{bj}R_4(g)(J_0)_{ij}\,dg.
\end{equation}
\begin{proposition}[Normalized stress tensor]
\label{ch14:normalized-stress}
The tensor $T$ is real, centered, symmetric, conserved and nonzero.
It transforms as a tensor under $R_4$. Its constant and linear local
fluxes are respectively the positional derivative and $L_d(K)$.
Along the same comparison sequence,
$U_{k,4}^{-1}(r_k^4T)\to T^0$.
\end{proposition}

\begin{proof}
Haar invariance gives the tensor law. All other linear properties
pass through the integral, conservation using vector covariance of
derivatives. In a transformed shell change variables by $g$ and use
correlation covariance; its actions are $\partial_b$ and $L_d(K)$.
Averaging preserves these normalized actions. If $T=0$, all tested
positional derivatives would vanish, contradicting smooth comparison
with a free differentiated quadratic field. Finally $T$ belongs to
the same centered conserved line as $J$. Its translation action
forces $J=c_*T$, and \eqref{ch14:balanced-stress} with
$c_k\to c_*$ proves the comparison limit.
\end{proof}

The complexified finite-dimensional rotation representations also
specify the label representations for the later continuation to
Lorentz transformations. No Lorentz action on a physical Hilbert
space, uniqueness of the vacuum, or positive extension of collision
kernels has been constructed in this chapter.
